\subsection*{Elegir un grupo equilibrado}
Para empezar, combina dos combatientes delante con un clérigo y un mago detrás. El clérigo aporta curación y comida; el mago, ataques y protecciones. Un ladrón o una combinación de clases permite añadir ganzúas sin renunciar necesariamente a otra función. Lleva armas a distancia para la retaguardia.

\begin{reference}{Clase}{Papel en el grupo}
Guerrero & Especialista en armas y armadura; adecuado para recibir los golpes de frente.\\
Explorador & Combatiente que comparte muchas opciones de equipo del guerrero. No presupongas que dispone del libro de un mago.\\
Paladín & Combatiente bueno, con Imposición de manos y acceso a su selección de oraciones. Su alineamiento limita la composición del grupo.\\
Clérigo & Curaciones, alimento, protecciones y ayuda contra muertos vivientes. Puede llevar armadura, pero su elección de armas es más limitada.\\
Mago & Aprende conjuros de su libro y de pergaminos. Sus armas y protección física son limitadas: protégelo en la retaguardia.\\
Ladrón & Especialista en cerraduras con ganzúas. No todas las cerraduras pueden abrirse así; conserva también las llaves.\\
Varias clases & Un personaje reúne funciones de las clases indicadas. La experiencia se reparte entre ellas y su progreso no equivale al de una sola clase.\\
\end{reference}

\textbf{Razas.} Humanos, elfos, semielfos, enanos, gnomos y medianos ofrecen distintas combinaciones de clases y límites de características. El menú de creación muestra las posibilidades reales de cada elección. Un elfo no puede recibir el conjuro Resucitar de esta versión; tenlo presente antes de crear el grupo. El sexo y el retrato permiten definir al personaje, sin cambiar cómo se controla.

\begin{reference}{Característica}{Qué debes mirar}
Fuerza (STR) & Importante para el combate físico y para forzar puertas. Algunos guerreros muestran además una fuerza excepcional junto al 18.\\
Inteligencia (INT) & Característica del mago; interesa al elegir un personaje destinado a estudiar magia.\\
Sabiduría (WIS) & Característica del clérigo; interesa para su capacidad de preparar oraciones.\\
Destreza (DEX) & Ayuda a la protección del personaje. Comprueba su efecto en la CA, sobre todo si lleva poca armadura.\\
Constitución (CON) & Importante para la resistencia y los puntos de vida. No descuides los PV de quienes van delante.\\
Carisma (CHA) & Parte de la ficha y de los requisitos de algunas clases. No es una orden para negociar ni sustituye las respuestas de los diálogos.\\
\end{reference}

\textbf{Alineamiento.} Describe la orientación moral del personaje y restringe algunas clases y combinaciones del grupo. No añade un botón de habilidad. Puedes revisar los personajes antes de pulsar JUGAR; no hace falta aceptar la primera tirada.
